\documentclass[11pt,a4paper]{article}

\usepackage[T1]{fontenc}
\usepackage[utf8]{inputenc}
\usepackage{lmodern}
\usepackage{amssymb,amsmath,amsthm,mathtools}
\usepackage[margin=2.5cm]{geometry}
\usepackage{parskip}
\usepackage[hidelinks]{hyperref}
\usepackage{microtype}

% Theorem environments (arXiv / journal style)
\theoremstyle{plain}
\newtheorem{theorem}{Theorem}[section]
\newtheorem{proposition}[theorem]{Proposition}
\newtheorem{lemma}[theorem]{Lemma}
\newtheorem{corollary}[theorem]{Corollary}

\theoremstyle{definition}
\newtheorem{definition}[theorem]{Definition}
\newtheorem{example}[theorem]{Example}
\newtheorem{remark}[theorem]{Remark}

% Notation
\newcommand{\Ext}{\mathrm{Ext}}
\newcommand{\Shape}{\mathrm{Shape}}
\newcommand{\Pos}{\mathrm{Pos}}
\newcommand{\Type}{\mathbf{Type}}
\newcommand{\id}{\mathrm{id}}
\newcommand{\Cont}{\mathbf{Cont}}
\newcommand{\Hom}{\mathrm{Hom}}

\title{The Category of Containers:\\
  Morphisms, Products, and Coproducts}
\author{Extracted from a machine-checked development}
\date{}

\begin{document}
\maketitle

\begin{abstract}
We equip containers and their morphisms with the structure of a category
$\Cont$. Morphisms consist of a forward map on shapes and a \emph{contravariant}
map on positions; with this variance, identity and composition laws hold
strictly. We then construct binary coproducts and binary products in $\Cont$,
exhibiting the universal properties in full (computation rules and uniqueness).
The proofs of all stated results have been machine-checked; we record only the
mathematical content.
\end{abstract}

\tableofcontents
\bigskip

%======================================================================
\section{Introduction}
%======================================================================

Containers---pairs $(S,P)$ consisting of a shape type and a family of position
types---classify a class of polynomial endofunctors on types via their
extensions. Morphisms of containers are maps on shapes together with
contravariant maps on positions; these induce natural transformations of
extensions, and conversely (the representation theorem).

The present note packages that morphism data into a category $\Cont$, and
develops its elementary algebra: binary coproducts and binary products. The
subtlety throughout is the variance of the position component. Because
positions are pulled back, the composite of two morphisms runs position maps
\emph{backwards}; once that convention is fixed, the category laws hold by
ordinary function associativity, with no transports.

\paragraph{What this note claims.}
In ordinary mathematical language: containers and shape-and-position morphisms
form a category; that category has binary coproducts (shapes a disjoint union,
positions case-wise) and binary products (shapes a cartesian product, positions
a \emph{coproduct} of the position sets). The universal properties are the
expected bijections
\[
  \Hom(C\oplus D,\,E)
  \;\cong\;
  \Hom(C,E)\times\Hom(D,E),
  \qquad
  \Hom(R,\,C\times D)
  \;\cong\;
  \Hom(R,C)\times\Hom(R,D).
\]

%======================================================================
\section{Containers and morphisms}
%======================================================================

We work in a single universe $\Type$. Recall:

\begin{definition}[Container]\label{def:container}
A \textbf{container} $C$ consists of a type $\Shape(C)$ of shapes and, for each
shape $s$, a type $\Pos(C)(s)$ of positions. We write $C=(S,P)$ when convenient.
\end{definition}

\begin{definition}[Container morphism]\label{def:morphism}
Let $C=(S_C,P_C)$ and $D=(S_D,P_D)$ be containers. A \textbf{morphism}
$\varphi:C\Rightarrow D$ consists of
\begin{enumerate}
  \item a map on shapes $u:S_C\to S_D$, and
  \item for each $s\in S_C$, a function
        $g_s:P_D\bigl(u(s)\bigr)\to P_C(s)$
        (contravariant on positions).
\end{enumerate}
We write $\varphi=(u,g)$, and refer to $u$ as the shape component and $g$ as
the position component.
\end{definition}

\begin{definition}[Identity morphism]\label{def:id}
For a container $C$, the \textbf{identity morphism} $\id_C:C\Rightarrow C$ has
shape component the identity function on $\Shape(C)$ and position component,
at each shape, the identity function on positions.
\end{definition}

\begin{definition}[Composition]\label{def:comp}
Let $\varphi=(u,g):C\Rightarrow D$ and $\psi=(v,h):D\Rightarrow E$ be
container morphisms. Their \textbf{composite}
$\psi\circ\varphi:C\Rightarrow E$ (written in applicative order; diagrammatic
order is $\varphi$ then $\psi$) has
\begin{enumerate}
  \item shape component $v\circ u$, and
  \item position component, at $s\in\Shape(C)$,
        \[
          g_s \;\circ\; h_{u(s)}
          \;:\;
          \Pos(E)\bigl(v(u(s))\bigr)
          \;\longrightarrow\;
          \Pos(C)(s).
        \]
\end{enumerate}
Shape maps compose forwards; position maps compose backwards.
\end{definition}

\begin{remark}
The backwards composition of position maps is forced by the intended
semantics: a labelling of $C$-positions is reindexed first by $\varphi$'s
pullback and then by $\psi$'s, or equivalently by the composite pullback.
\end{remark}

Because the type of the position component depends on the shape component, an
extensionality principle for morphisms must account for that dependence.

\begin{lemma}[Extensionality for morphisms]\label{lem:ext}
Let $\varphi,\psi:C\Rightarrow D$ be container morphisms. If the shape
components agree and the position components are (heterogeneously) equal, then
$\varphi=\psi$.
\end{lemma}

\begin{proof}
By lean.
\end{proof}

\begin{lemma}[Pointwise extensionality]\label{lem:ext-prime}
Let $\varphi,\psi:C\Rightarrow D$ with shape components $u$ and $v$, and
suppose $h:u=v$. If for every shape $s$ and every $D$-position $p$ of $v(s)$,
\[
  g_s(p) \;=\; g'_s\bigl((h\cdot s)\cdot p\bigr)
\]
(after transporting $p$ along the action of $h$ at $s$), then $\varphi=\psi$.
\end{lemma}

\begin{proof}
By lean.
\end{proof}

%======================================================================
\section{The category of containers}
%======================================================================

\begin{proposition}[Left unit]\label{prop:id-comp}
For every container morphism $\varphi:C\Rightarrow D$,
\begin{equation}
  \varphi \;\circ\; \id_C \;=\; \varphi.
\end{equation}
\end{proposition}

\begin{proof}
By lean.
\end{proof}

\begin{proposition}[Right unit]\label{prop:comp-id}
For every container morphism $\varphi:C\Rightarrow D$,
\begin{equation}
  \id_D \;\circ\; \varphi \;=\; \varphi.
\end{equation}
\end{proposition}

\begin{proof}
By lean.
\end{proof}

\begin{proposition}[Associativity]\label{prop:assoc}
For morphisms $\varphi:C\Rightarrow D$, $\psi:D\Rightarrow E$,
$\chi:E\Rightarrow F$,
\begin{equation}
  \chi \;\circ\; (\psi \;\circ\; \varphi)
  \;=\;
  (\chi \;\circ\; \psi) \;\circ\; \varphi.
\end{equation}
\end{proposition}

\begin{proof}
By lean.
\end{proof}

\begin{remark}
All three laws hold strictly: the underlying shape-map operations are ordinary
function identity and composition, which are unital and associative on the
nose, and the backwards position maps inherit the same equalities with no
transport.
\end{remark}

\begin{definition}[The category $\Cont$]\label{def:cont}
The category $\Cont$ has
\begin{enumerate}
  \item objects: containers;
  \item morphisms $C\to D$: container morphisms $C\Rightarrow D$;
  \item identities and composition as in Definitions~\ref{def:id}
        and~\ref{def:comp};
  \item unit and associativity laws as in Propositions~\ref{prop:id-comp},
        \ref{prop:comp-id}, and~\ref{prop:assoc}.
\end{enumerate}
\end{definition}

%======================================================================
\section{Binary coproducts}
%======================================================================

\begin{definition}[Coproduct of containers]\label{def:coprod}
Let $C=(S,P)$ and $D=(T,Q)$ be containers. Their \textbf{coproduct}
$C\oplus D$ is the container with
\begin{enumerate}
  \item shapes $S+T$ (the disjoint union), and
  \item positions given case-wise:
        $\Pos(C\oplus D)(\mathsf{inl}\,s)=P(s)$ and
        $\Pos(C\oplus D)(\mathsf{inr}\,t)=Q(t)$.
\end{enumerate}
\end{definition}

\begin{definition}[Injections]\label{def:inl-inr}
The \textbf{left injection} $\iota_C:C\Rightarrow C\oplus D$ has shape
component $\mathsf{inl}:S\to S+T$ and position component the identity on
$P(s)$ (positions of an $\mathsf{inl}$-shape \emph{are} the $C$-positions).
The \textbf{right injection} $\iota_D:D\Rightarrow C\oplus D$ is defined
dually via $\mathsf{inr}$.
\end{definition}

\begin{definition}[Copairing]\label{def:desc}
Given morphisms $\varphi:C\Rightarrow E$ and $\psi:D\Rightarrow E$, their
\textbf{copairing} $[\varphi,\psi]:C\oplus D\Rightarrow E$ has shape
component $\mathsf{elim}(u_\varphi,u_\psi)$ (dispatch on the sum) and
position component that of $\varphi$ on $\mathsf{inl}$-shapes and that of
$\psi$ on $\mathsf{inr}$-shapes.
\end{definition}

\begin{proposition}[Left computation]\label{prop:inl-desc}
For morphisms $\varphi:C\Rightarrow E$ and $\psi:D\Rightarrow E$,
\begin{equation}
  [\varphi,\psi] \;\circ\; \iota_C \;=\; \varphi.
\end{equation}
\end{proposition}

\begin{proof}
By lean.
\end{proof}

\begin{proposition}[Right computation]\label{prop:inr-desc}
For morphisms $\varphi:C\Rightarrow E$ and $\psi:D\Rightarrow E$,
\begin{equation}
  [\varphi,\psi] \;\circ\; \iota_D \;=\; \psi.
\end{equation}
\end{proposition}

\begin{proof}
By lean.
\end{proof}

\begin{proposition}[Uniqueness / $\eta$-rule]\label{prop:desc-eta}
For every morphism $\chi:C\oplus D\Rightarrow E$,
\begin{equation}
  \bigl[\,\chi\circ\iota_C,\; \chi\circ\iota_D\,\bigr] \;=\; \chi.
\end{equation}
\end{proposition}

\begin{proof}
By lean.
\end{proof}

\begin{theorem}[Universal property of the coproduct]\label{thm:coprod-up}
For every cocone $(\varphi:C\Rightarrow E,\;\psi:D\Rightarrow E)$ there is a
unique morphism $\chi:C\oplus D\Rightarrow E$ such that
$\chi\circ\iota_C=\varphi$ and $\chi\circ\iota_D=\psi$. Explicitly,
$\chi=[\varphi,\psi]$; uniqueness follows from Propositions~\ref{prop:inl-desc},
\ref{prop:inr-desc}, and~\ref{prop:desc-eta}. Equivalently,
\begin{equation}
  \Hom_{\Cont}(C\oplus D,\,E)
  \;\cong\;
  \Hom_{\Cont}(C,E)\times\Hom_{\Cont}(D,E)
\end{equation}
naturally in $E$.
\end{theorem}

\begin{proof}
By lean.
\end{proof}

%======================================================================
\section{Binary products}
%======================================================================

The product of containers is dual in the categorical sense, but the
combinatorics of positions are not a naive product: the positions of a
product shape are the \emph{coproduct} of the two position sets. This is the
product of the underlying polynomials.

\begin{definition}[Product of containers]\label{def:prod}
Let $C=(S,P)$ and $D=(T,Q)$ be containers. Their \textbf{product}
$C\times D$ is the container with
\begin{enumerate}
  \item shapes $S\times T$, and
  \item positions of $(s,t)$ given by $P(s)+Q(t)$.
\end{enumerate}
\end{definition}

\begin{definition}[Projections]\label{def:fst-snd}
The \textbf{first projection} $\pi_C:C\times D\Rightarrow C$ has shape
component the first projection $S\times T\to S$, and position component
(contravariant) the coproduct injection $\mathsf{inl}:P(s)\to P(s)+Q(t)$.
The \textbf{second projection} $\pi_D$ is defined dually via $\mathsf{inr}$.
\end{definition}

\begin{remark}
The variance is the teaching point dual to the coproduct. There the injections
were identity-on-positions; here the projections carry a coproduct injection
\emph{backwards} on positions: a $C$-position is sent into the larger
position set $P(s)+Q(t)$.
\end{remark}

\begin{definition}[Pairing]\label{def:pair}
Given morphisms $f:R\Rightarrow C$ and $g:R\Rightarrow D$, their
\textbf{pairing} $\langle f,g\rangle:R\Rightarrow C\times D$ has shape
component $r\mapsto\bigl(u_f(r),\,u_g(r)\bigr)$ and position component, at
$r$, the coproduct elimination of $f$'s and $g$'s position maps.
\end{definition}

\begin{proposition}[Left computation]\label{prop:pair-fst}
For morphisms $f:R\Rightarrow C$ and $g:R\Rightarrow D$,
\begin{equation}
  \pi_C \;\circ\; \langle f,g\rangle \;=\; f.
\end{equation}
\end{proposition}

\begin{proof}
By lean.
\end{proof}

\begin{proposition}[Right computation]\label{prop:pair-snd}
For morphisms $f:R\Rightarrow C$ and $g:R\Rightarrow D$,
\begin{equation}
  \pi_D \;\circ\; \langle f,g\rangle \;=\; g.
\end{equation}
\end{proposition}

\begin{proof}
By lean.
\end{proof}

\begin{proposition}[Uniqueness / $\eta$-rule]\label{prop:pair-eta}
For every morphism $\chi:R\Rightarrow C\times D$,
\begin{equation}
  \bigl\langle\,
    \pi_C\circ\chi,\;
    \pi_D\circ\chi
  \,\bigr\rangle
  \;=\; \chi.
\end{equation}
\end{proposition}

\begin{proof}
By lean.
\end{proof}

\begin{theorem}[Universal property of the product]\label{thm:prod-up}
For every cone $(f:R\Rightarrow C,\;g:R\Rightarrow D)$ there is a unique
morphism $\chi:R\Rightarrow C\times D$ such that $\pi_C\circ\chi=f$ and
$\pi_D\circ\chi=g$. Explicitly, $\chi=\langle f,g\rangle$; uniqueness follows
from Propositions~\ref{prop:pair-fst}, \ref{prop:pair-snd},
and~\ref{prop:pair-eta}. Equivalently,
\begin{equation}
  \Hom_{\Cont}(R,\,C\times D)
  \;\cong\;
  \Hom_{\Cont}(R,C)\times\Hom_{\Cont}(R,D)
\end{equation}
naturally in $R$.
\end{theorem}

\begin{proof}
By lean.
\end{proof}

%======================================================================
\section{Discussion}
%======================================================================

\paragraph{What has been established.}
Containers and their morphisms form a category $\Cont$ with strictly holding
unit and associativity laws. Binary coproducts and products exist and satisfy
the full universal properties. The combinatorics are those of polynomial
functors: coproduct of containers is case-wise on both shapes and positions;
product of containers multiplies shapes and \emph{adds} positions.

\paragraph{Variance.}
The entire development turns on the contravariance of the position component.
With positions pulled back, composites and universal maps typecheck without
transport, and all category and universal-property laws reduce to function
extensionality and case analysis on sum types.

\paragraph{Scope.}
This note does not treat the composition product of containers (the monoidal
structure whose comonoids are directed containers), nor the extension functor
$\Cont\to[\Type,\Type]$ as a fully faithful embedding. Those belong to
companion developments.

\paragraph{Machine checking.}
Every definition and theorem statement above corresponds to a declaration in a
machine-checked formalisation. The phrase ``by lean'' indicates that the proof
has been verified by a proof assistant.

%======================================================================
\section*{Acknowledgements}
%======================================================================

The mathematical content follows the theory of containers of Abbott,
Altenkirch, and Ghani. The presentation here is a prose rendering of a formal
development intended for human review of the theorem statements.

\begin{thebibliography}{9}

\bibitem{abbott2003}
M.~Abbott, T.~Altenkirch, and N.~Ghani.
\newblock Categories of containers.
\newblock In \emph{FoSSaCS}, volume~2620 of \emph{LNCS}, pages 23--38.
\newblock Springer, 2003.

\bibitem{abbott2005}
M.~Abbott, T.~Altenkirch, and N.~Ghani.
\newblock Containers: constructing strictly positive types.
\newblock \emph{Theoretical Computer Science}, 342(1):3--27, 2005.

\end{thebibliography}

\end{document}
